\textbf{Phần 1. Trắc nghiệm lựa chọn}
\setcounterpageref{ex}{0}

\begin{ex} %[3D1N01]
Số \textit{Bốn trăm năm mươi sáu} viết là
\choice
{$465$}
{\True $456$}
{$546$}
{$4056$}
\loigiai{
Số bốn trăm năm mươi sáu viết là $456$.\\
\textbf{Đáp án của câu hỏi là B}
}
\end{ex}

\begin{ex} %[3D1N01]
Số liền sau của số $399$ là
\choice
{\True $400$}
{$398$}
{$401$}
{$390$}
\loigiai{
Số liền sau của $399$ là $399 + 1 = 400$.\\
\textbf{Đáp án của câu hỏi là A}
}
\end{ex}

\begin{ex} %[3D1H01]
Trong số $782$, chữ số 8 có giá trị là
\choice
{$8$}
{$800$}
{$8000$}
{\True $80$}
\loigiai{
Chữ số 8 ở hàng chục nên có giá trị là $80$.\\
\textbf{Đáp án của câu hỏi là D}
}
\end{ex}

\begin{ex} %[3D1H01]
Số gồm $5$ trăm, $0$ chục và $7$ đơn vị viết là
\choice
{$570$}
{$57$}
{\True $507$}
{$705$}
\loigiai{
Số đó là $507$.\\
\textbf{Đáp án của câu hỏi là C}
}
\end{ex}

\begin{ex} %[3D1H01]
Điền dấu thích hợp vào chỗ chấm: $685 \dots 658$.
\choice
{\True $>$}
{$<$}
{$=$}
{Không so sánh được}
\loigiai{
Hàng trăm bằng nhau, so sánh hàng chục: $8 > 5$ nên $685 > 658$.\\
\textbf{Đáp án của câu hỏi là A}
}
\end{ex}

\begin{ex} %[3D1H01]
Số lớn nhất có 3 chữ số là
\choice
{$100$}
{$987$}
{\True $999$}
{$1000$}
\loigiai{
Số lớn nhất có 3 chữ số là $999$.\\
\textbf{Đáp án của câu hỏi là C}
}
\end{ex}

\begin{ex} %[3D1V01]
Sắp xếp các số $452, 542, 245, 524$ theo thứ tự từ bé đến lớn:
\choice
{$542 < 524 < 452 < 245$}
{$245 < 452 < 542 < 524$}
{$452 < 245 < 524 < 542$}
{\True $245 < 452 < 524 < 542$}
\loigiai{
Thứ tự từ bé đến lớn: $245 < 452 < 524 < 542$.\\
\textbf{Đáp án của câu hỏi là D}
}
\end{ex}

\begin{ex} %[3D1V01]
Số bé nhất có 3 chữ số khác nhau là
\choice
{$100$}
{\True $102$}
{$123$}
{$101$}
\loigiai{
Số bé nhất có 3 chữ số khác nhau là $102$.\\
\textbf{Đáp án của câu hỏi là B}
}
\end{ex}

\begin{ex} %[3D1V01]
Tìm số tròn chục lớn nhất có 3 chữ số:
\choice
{$999$}
{$900$}
{\True $990$}
{$980$}
\loigiai{
Số tròn chục lớn nhất có 3 chữ số là $990$.\\
\textbf{Đáp án của câu hỏi là C}
}
\end{ex}

\begin{ex} %[3D1V01]
Từ ba chữ số $3, 0, 7$ lập được bao nhiêu số có 3 chữ số khác nhau?
\choice
{$6$ số}
{$3$ số}
{$5$ số}
{\True $4$ số}
\loigiai{
Chữ số hàng trăm có 2 cách chọn (3 hoặc 7). Các số lập được là: $307, 370, 703, 730$ (4 số).\\
\textbf{Đáp án của câu hỏi là D}
}
\end{ex}

\begin{ex} %[3D1V01]
Dãy số nào sau đây được viết theo thứ tự giảm dần?
\choice
{$312; 321; 330; 340$}
{\True $876; 867; 786; 678$}
{$123; 234; 345; 456$}
{$500; 550; 505; 501$}
\loigiai{
Dãy giảm dần: $876 > 867 > 786 > 678$.\\
\textbf{Đáp án của câu hỏi là B}
}
\end{ex}

\begin{ex} %[3D1C01]
Tìm số có 3 chữ số biết chữ số hàng trăm là số lớn nhất có 1 chữ số, chữ số hàng chục gấp đôi chữ số hàng đơn vị và chữ số hàng đơn vị là 3.
\choice
{\True $963$}
{$936$}
{$863$}
{$962$}
\loigiai{
Hàng trăm là 9. Hàng đơn vị là 3. Hàng chục là $3 \times 2 = 6$. Số đó là $963$.\\
\textbf{Đáp án của câu hỏi là A}
}
\end{ex}

\textbf{Phần 2. Trắc nghiệm đúng sai}
\setcounterpageref{ex}{0}

\begin{ex} %[3D1H01]
Xét tính đúng sai của các nhận định về số có ba chữ số:
\choiceTF
{\True Số $100$ là số bé nhất có 3 chữ số.}
{\True Số liền trước của số $500$ là số $499$.}
{Số $730$ gồm 7 trăm và 3 đơn vị.}
{\True Số $730$ gồm 7 trăm và 3 chục.}
\loigiai{
\begin{itemize}
    \item $100$ là số bé nhất có 3 chữ số: Đúng.
    \item $500 - 1 = 499$: Đúng.
    \item Số $730$ gồm 7 trăm và 3 chục (không phải 3 đơn vị). Ý 3 Sai, Ý 4 Đúng.
\end{itemize}
\begin{itemize}
    \item a) \textbf{Đúng}.
    \item b) \textbf{Đúng}.
    \item c) \textbf{Sai}.
    \item d) \textbf{Đúng}.
\end{itemize}
}
\end{ex}

\begin{ex} %[3D1V01]
Cho 4 số: $345, 543, 435, 354$. Xét tính đúng sai:
\choiceTF
{\True Số lớn nhất trong 4 số là $543$.}
{\True Số bé nhất trong 4 số là $345$.}
{Có 3 số lớn hơn $400$.}
{\True Có 2 số lớn hơn $400$ (là $435$ và $543$).}
\loigiai{
\begin{itemize}
    \item Số lớn nhất là 543: Đúng.
    \item Số bé nhất là 345: Đúng.
    \item Các số lớn hơn 400 là 435 và 543 (có 2 số). Ý 3 Sai, Ý 4 Đúng.
\end{itemize}
\begin{itemize}
    \item a) \textbf{Đúng}.
    \item b) \textbf{Đúng}.
    \item c) \textbf{Sai}.
    \item d) \textbf{Đúng}.
\end{itemize}
}
\end{ex}

\textbf{Phần 3. Trả lời ngắn}
\setcounterpageref{ex}{0}

\begin{ex} %[3D1H01]
Viết số liền trước của số $800$.
\shortans{799}
\loigiai{
\textbf{Đáp số:} $799$\\
$800 - 1 = 799$.
}
\end{ex}

\begin{ex} %[3D1H01]
Viết số gồm $6$ trăm và $5$ đơn vị.
\shortans{605}
\loigiai{
\textbf{Đáp số:} $605$\\
$605$.
}
\end{ex}

\begin{ex} %[3D1V01]
Tìm số lớn nhất có 3 chữ số khác nhau được lập từ các chữ số $2, 8, 5$.
\shortans{852}
\loigiai{
\textbf{Đáp số:} $852$\\
$852$.
}
\end{ex}

\begin{ex} %[3D1V01]
Tìm số bé nhất có 3 chữ số khác nhau được lập từ các chữ số $6, 0, 9$.
\shortans{609}
\loigiai{
\textbf{Đáp số:} $609$\\
$609$.
}
\end{ex}

\begin{ex} %[3D1V01]
Dãy số $100, 102, 104, \dots, 120$ có bao nhiêu số hạng?
\shortans{11}
\loigiai{
\textbf{Đáp số:} $11$\\
$(120 - 100) : 2 + 1 = 11$ số.
}
\end{ex}

\begin{ex} %[3D1C01]
Tìm số có 3 chữ số, biết tổng ba chữ số bằng 3 và chữ số hàng trăm lớn hơn chữ số hàng chục, chữ số hàng chục lớn hơn chữ số hàng đơn vị.
\shortans{210}
\loigiai{
\textbf{Đáp số:} $210$\\
Ta có $2 + 1 + 0 = 3$ và $2 > 1 > 0$. Số đó là $210$.
}
\end{ex}
